\begin{abstract}
Archaic hominin introgression has shaped the genome of modern humans,
yet existing detection methods quantify introgression through archaic variant density,
discarding the haplotype match-length signal that coalescent theory predicts to
differ $\sim$11-fold between introgressed and non-introgressed segments.
Here we introduce \textbf{ArchaicPainter} (\AP{}), a three-state
(AMH / Neanderthal / Denisovan) Hidden Markov Model that adapts the
tract-length transition structure of the Li \& Stephens haplotype copying model
to archaic ancestry inference, with hidden states representing ancestry
categories rather than individual reference haplotypes.
In 50 independent simulation replicates with known ground truth,
\AP{} achieves segment-level F1 $= 0.480 \pm 0.095$ (AUPRC $= 0.577$)
(AUPRC $= 0.577$),
Ablation analysis reveals that the gain derives from a bidirectional emission model
exploiting both matches and mismatches as archaic-state evidence;
a positive-only emission mode --- appropriate when the available archaic reference
diverges from the true introgressing population --- is provided for real-data use.
Applied to chromosome~21 from 1000~Genomes Project Phase~3 data,
\AP{} detects $1.10\% \pm 0.60\%$ Neanderthal ancestry in Europeans (CEU)
and $0.98\% \pm 0.52\%$ in East Asians (CHB)
(median segment length $\sim$28--57~kb), consistent with published genome-wide
estimates and complementary to IBDmix's high-sensitivity short-tract detection.
A three-state extension achieves proof-of-concept simultaneous Neanderthal and
Denisovan source attribution (NEA F1 $= 0.260$, DEN F1 $= 0.188$),
and near-linear computational scaling (empirical exponent 1.21, $0.11$~s per
haplotype at $n = 1{,}000$) supports population-scale deployment.
As an illustrative application, functional annotation of detected chr21 regions
recovers 7 of 16 published adaptive introgression candidates and a complete
interferon receptor gene cluster
(\textit{IFNAR1}, \textit{IFNAR2}, \textit{IFNGR2}, \textit{IL10RB})
in an unsupervised manner, consistent with previously reported archaic
contributions to antiviral immunity; formal permutation-based enrichment
tests are needed to establish statistical significance.
\AP{} is released as open-source software with reproducible pipelines
for community reuse.
\end{abstract}
